%\documentclass[11pt,a4paper,twoside]{article}

\usepackage[utf8]{inputenc}

\usepackage[T1,T2A]{fontenc}

\usepackage[english.ancient,russian.ancient]{babel}

\usepackage{graphicx}

\usepackage{array}

\usepackage[doi]{samgtu-bib}

\usepackage{samgtu}

\usepackage{samgtu-name}

\usepackage{mathtools}

\mathtoolsset{showonlyrefs}

\usepackage{desclist}

\usepackage{euscript}

\usepackage{mathrsfs} 

\usepackage{multicol}

\usepackage{mathrsfs}

\usepackage{cite}

\usepackage{cmap}

\usepackage{qrcode}

\allowdisplaybreaks[4]

\usepackage[nodayofweek]{datetime}

\usepackage[thinlines]{easybmat}



\renewcommand{\JournalYear}{2018}

\renewcommand{\JournalVolume}{22}

\renewcommand{\JournalNumber}{4}



\newdate{JournalDateSign}{29}{12}{2018}% Дата подписания Вестника в печать

\renewcommand{\JournalDateSign}{\displaydate{JournalDateSign}}

\newdate{JournalDatePrint}{16}{01}{2019}% Дата выхода Вестника из печати

\renewcommand{\JournalDatePrint}{\displaydate{JournalDatePrint}}

%\renewcommand{\JournalUPL}{16.56} % Усл. печ. л.

%\renewcommand{\JournalUIL}{16.53} % Уч.-изд. л.

\renewcommand{\JournalUPL}{15.24} % Усл. печ. л.

\renewcommand{\JournalUIL}{15.21} % Уч.-изд. л.

\renewcommand{\JournalRegNumber}{220/18} % Регистрационный номер

\renewcommand{\JournalOrderNumber}{2} % Номер заказа



%\renewcommand{\JournalRMonth}{Март} % Месяц выхода Вестника

%\renewcommand{\JournalEMonth}{March} % Месяц выхода Вестника

%\renewcommand{\JournalRMonth}{Июнь} % Месяц выхода Вестника

%\renewcommand{\JournalEMonth}{June} % Месяц выхода Вестника

%\renewcommand{\JournalRMonth}{Сентябрь} % Месяц выхода Вестника

%\renewcommand{\JournalEMonth}{September} % Месяц выхода Вестника

\renewcommand{\JournalRMonth}{Декабрь} % Месяц выхода Вестника

\renewcommand{\JournalEMonth}{December} % Месяц выхода Вестника



\newcommand{\totop}[1]{\raisebox{6pt}[1pt][2pt]{#1}}

\newcommand{\tottop}[1]{\raisebox{12pt}[1pt][2pt]{#1}} 

\newcommand{\totttop}[1]{\raisebox{18pt}[1pt][2pt]{#1}} 

\newcommand{\tottttop}[1]{\raisebox{24pt}[1pt][2pt]{#1}} 

\newcommand{\totttttop}[1]{\raisebox{30pt}[1pt][2pt]{#1}} 

\newcommand{\tobot}[1]{\raisebox{-6pt}[1pt][2pt]{#1}} 

\newcommand{\tobbot}[1]{\raisebox{-12pt}[1pt][2pt]{#1}} 

\newcommand{\tobbbot}[1]{\raisebox{-18pt}[1pt][2pt]{#1}}



\graphicspath{{./00PIC/}}

%

\vsgtu{1639} %registration number

\FirstPage{774}

\LastPage{784}

%

%\usepackage{samgtu-name}

%

%\pdfcrop

%\draft

%

\ShortCommunication

%

%\begin{document}





$\;$

\vspace{1mm}



\hfill \qrcode[hyperlink,height=.8in]{https://doi.org/10.14498/vsgtu1639}

\vspace{-.8in}



\udc{517.955, 517.968.7}

\msc{26A33, 35K15, 35R11}



\rutitle{К проблеме единственности решения задачи Коши для~уравнения дробной диффузии с оператором Бесселя}

\rutitleshort{К проблеме единственности решения задачи Коши\dots}

\rutitlehack{К проблеме единственности решения \\ задачи Коши для~уравнения дробной диффузии с~оператором Бесселя}



\entitle{On the uniqueness of the solution of the Cauchy problem for the equation of fractional diffusion with~Bessel~operator}

\entitleshort{On the uniqueness of the solution of the Cauchy problem\dots}



\ruauthor{Ф.~Г.~Хуштова}{Х\,у\,ш\,т\,о\,в\,а~Ф.~Г.}

\enauthor{F.~G.~Khushtova}{K\,h\,u\,s\,h\,t\,o\,v\,a~F.~G.}



\ruaffil{\runiipma.}



\enaffil{\enniipma.}



\ruauthorsdetails{1}{% Количество авторов

\fullauthor{\rupid{53181}{Фатима Гидовна Хуштова}}

\CAuthor % Автор-корреспондент

\orcid{0000-0003-4088-3621} % ORCID ID автора 

%\regalia{без ученой степени} 

\position[n]{научный сотрудник} 

\dept{отдел дробного исчисления} 

\email{khushtova@yandex.ru}

}



\enauthorsdetails{1}{

\fullauthor{\enpid{53181}{Fatima G. Khushtova}}

\CAuthor % Сorresponding Author

\orcid{0000-0003-4088-3621} % ORCID ID

%\regalia{Cand. Phys. \& Math. Sci., Associate Professor} 

\position[n]{Researcher} 

\dept{Dept. of Fractional Calculus}

\email{khushtova@yandex.ru}

}



\rukeywords{уравнение дробной диффузии, оператор дробного дифференцирования, оператор Бесселя, задача Коши, единственность решения, условие Тихонова, функция Райта}

\enkeywords{fractional diffusion equation, fractional differentiation operator, Bessel operator, Cauchy problem, solution uniqueness, Tikhonov condition,	 Wright function}



\ruabstract{Рассматривается уравнение дробной диффузии с сингулярным оператором Бесселя, действующим по пространственной переменной, и оператором дробного дифференцирования Римана--Лиувилля, действующим по временной переменной. 

Когда порядок дробной производной равен единице, а особенность у оператора Бесселя отсутствует, рассматриваемое уравнение совпадает с~классическим уравнением теплопроводности. 

Ранее для уравнения дробной диффузии с оператором Бесселя было построено решение задачи Коши и доказана теорема единственности решения в классе функций экспоненциального роста. 



Построен пример, показывающий, что увеличение показателя степени в условии, гарантирующем единственность решения задачи Коши, влечет за собой неединственность решения. 

С помощью известных свойств функции Райта получены оценки для построенной функции.

Показывается, что она, будучи не равной тождественно нулю, удовлетворяет однородному уравнению и однородному условию Коши. 

}





\enabstract{In this paper, we consider fractional diffusion equation involving the Bessel operator acting with respect to a spatial variable and the Riemann-Liouville fractional differentiation operator acting with respect to a time variable. When the order of the fractional derivative is unity, and the singularity of the Bessel operator is absent, this equation coincides with the classical heat equation. Earlier, a solution of the Cauchy problem has been considered for the considered equation and a uniqueness theorem has been proved for a class of functions satisfying the analog of the Tikhonov condition. 



In this paper, we have constructed an example to show that the exponent (power) at the condition of the uniqueness of the solution to the Cauchy problem cannot be raised under. Its increase leads to a non-uniqueness of the solution. Using the well-known properties of the Wright function, we have obtained estimates for constructed function, which satisfies the homogeneous equation and the zero Cauchy condition.

}



\newdate{khusha}{28}{08}{2018}

\date{\displaydate{khusha}}

\newdate{khushb}{25}{10}{2018}

\revisiondate{\displaydate{khushb}}

\newdate{khushc}{12}{11}{2018}

\accepteddate{\displaydate{khushc}}



\newdate{khushe}{28}{11}{2018}

\renewcommand{\JournalDataOnline}{\displaydate{khushe}}



%\ForPeerReview

\makerutitle



\newpage



\Section[n]{Введение}

В области $\Omega=\{(x,y): -\infty<x<\infty,\ 0<y<T \}$ рассмотрим уравнение

\begin{equation}\label{GL3_1}

B_xu(x,y)-D^{\alpha}_{0y}u(x,y)=0, 

\end{equation}

где $D_{0y}^{\alpha}$ "--- оператор дробного дифференцирования в смысле  

Римана"--~Лиувилля  порядка $\alpha,$ определяемый равенствами~\cite{Drob, Pskhu}   

$$D_{0y}^{\alpha} \, g(y)= \frac{dg}{dy},$$ если $\alpha=1$, и

$$ 

D_{0y}^{\alpha}\,g(y)=\frac{1}{\Gamma(1-\alpha)}\frac{d}{dy}

\int _{0}^{y}\frac{g(t)}{(y-t)^{\alpha}}dt,

$$ 

если $0<\alpha<1;$

$$

B_x=x^{-b} \frac{d}{dx}\left(x^b \frac{d}{dx}\right)=\frac{d^2}{d x^2}+\frac{b}{x} \frac{d}{d x}

$$

"--- оператор Бесселя, $|b|<1.$



Уравнение \eqref{GL3_1} при $\alpha=1,$ то есть уравнение

\begin{equation}\label{UrAlpha-1}

u_{xx}(x,y)+\frac{b}{x}u_{x}(x,y)-u_{y}(x,y)=0,

\end{equation}

совпадает с %~названным И.~А.~Киприяновым 

$B$-параболическим уравнением.

Краевые задачи в ограниченной и~неограниченной областях для этого уравнения при различных значениях параметра $b$ рассматривали многие авторы. 

Например, в работе V.~Alexia\-des~\cite{Alexiades} для него решены первая, вторая и~третья краевые задачи в~области с~подвижной границей, в~работе D.~Calton\cite{Calton} для него  исследована задача Коши.

 

 Уравнение \eqref{UrAlpha-1} заменой $z=x^2/(4n)$ сводится к~уравнению

\begin{equation} \label{UrKep2}

u_{zz}(z,y)+\frac{m+1}{z}\,u_{z}(z,y)-\frac{n}{z} u_{y}(z,y)=0,\quad m=(b-1)/2,

\end{equation}

для которого S.~K\c{e}pi\'nski  \cite{Kepinski_2} исследовал краевую задачу в~полуполосе $x>0,$ $0<y<y_0.$

Уравнение \eqref{UrAlpha-1} было также объектом исследования  монографий С.~А.~Терсенова~\cite{Tersenov} и~М.~И.~Матийчука~\cite{Matiychuk}.



Дифференциальные уравнения, содержащие оператор Бесселя, наиболее подробно исследованы в~работах И.~А.~Киприянова и~его учеников \cite{Kipriyanov-Monographiya,  Kipriyanov}.

Отметим здесь, что очень важные результаты по применению операторов преобразования в~теории дифференциальных уравнений с~особенностями в~коэффициентах, в~частности с~операторами Бесселя, получены В.~В.~Катраховым и~С.~М.~Ситником~\cite{KatrakhovSitnik, Sitnik}.





В случае, когда $b=0,$ $0<\alpha<2,$ уравнение \eqref{GL3_1} совпадает с диффузионно"=волновым уравнением. 

Различные краевые задачи для него, а также для многомерных его обобщений подробно исследованы в работах многих авторов.  Приведем некоторые из них. 



А.~Н.~Кочубей \cite{Kochubei_2} исследовал задачу Коши для уравнения диффузии дробного порядка с~регуляризованной дробной производной (производной Капуто) и~эллиптическим оператором с~непрерывными ограниченными вещественными коэффициентами

$$

L=\sum\limits_{i,j=1}^{n}a_{i,j}(x)\frac{\partial^2}{\partial x_i \partial x_j}+

\sum\limits_{j=1}^{n}b_{j}(x)\frac{\partial}{\partial x_j}+c(x) u.

$$

В случае, когда $L=\Delta$ "--- оператор Лапласа, фундаментальное решение построено в терминах $H$-функции Фокса. 

Исследованию задачи Коши для диффузионного и диффузионно"=волнового уравнений дробного порядка с~производной Капуто  посвящены также  работы А.~Н.~Кочубея и С.~Д.~Эйдельмана \cite{Kochubei_1, Kochubei_3, Kochubei_4, Kochubei_5}.



А.~В.~Псху~\cite{Pskhu2} построил фундаментальное решение и исследовал задачу Коши  для многомерного диффузионно"=волнового уравнения с~оператором дробного дифференцирования Джрбашяна"--~Нерсесяна, который  включает в~себя как частный случай операторы дробного дифференцирования Римана"--~Лиувилля и Капуто.



В~работах  \cite{Pskhu-SibirJournal, Pskhu3} также построено фундаментальное решение и исследована задача Коши для уравнения дробной диффузии соответственно с~оператором дискретно распределенного дифференцирования и оператором Джрбашяна"--~Нерсесяна, действующим по многим переменным времени.

В работе \cite{Pskhu4} решена первая краевая задача в~нецилиндрической области для диффузионно"=волнового уравнения с~оператором Джрбашяна"--~Нерсесяна. 

Доказана теорема существования и~единственности исследуемой задачи, конструктивно построено ее решение. 



А.~А.~Ворошилов и А.~А.~Килбас \cite{VorKilbas} методом интегральных преобразований построили решение задачи Коши для уравнения

\begin{equation} \label{MnogPerem}

D^{\alpha}_{0t}\,u(x,t)=\lambda^2 \Delta_x u(x,t), \quad  x\in{\mathbb R}^m,\quad  t>0, 

\end{equation}

где $\Delta_x=\sum_{j=1}^{m}\partial^2/\partial x_{j}^{2},$ $n-1<\alpha <n,$ $n\in\mathbb{N}.$ 

При $0<\alpha\leqslant 1$ и $1<\alpha<2$ решения  выписаны в терминах $H$-функции Фокса.

В работе~\cite{VorKilbas2} было также  построено решение задачи Коши в~случае, когда в уравнении

\eqref{MnogPerem} вместо оператора Римана"--~Лиувилля содержится оператор Капуто. 



Задача Коши для уравнения \eqref{MnogPerem} при $\lambda=1,$ 

$0<\alpha\leqslant 1$  была ранее рассмотрена С.~Х.~Геккиевой~\cite{Gekk}, а~в~работе \cite{Gekkieva1} ею исследована первая краевая задача в~первом квадранте для уравнения \eqref{MnogPerem} при $m=\lambda=1,$ 

$0<\alpha\leqslant 1.$  







Уравнение диффузии дробного порядка и некоторые его обобщения рассматривали 

F.~Mainardi, G.~Pagnini, Yu.~Luchko, P.~Paradisi~\cite{Mainardi-1, Mainardi-3, Mainardi-2, Pagnini_1,   Pagnini_Paradisi}  и~многие другие. 

Более подробную библиографию по данному вопросу можно найти в~монографиях \cite{Pskhu, MathaiSaxenaHaubold, Podlubny}.

 

%Чтобы подготовить Вашу статью для рецензента, т.~е. обезличить её, воспользуйтесь командой \verb"\ForPeerReview" перед командой 

%\verb"\makerutitle".



\smallskip





\Section{Постановка задачи Коши и теорема единственности решения}

Пусть $\Omega^{+}=\Omega\cap\{x>0 \},$ $\Omega^{-}=\Omega\cap\{x<0\},$ 

$\bar{\Omega}$ "--- замыкание области $\Omega.$ 



\smallskip



{\sc\small Определение.}

{\it Регулярным  решением} уравнения \eqref{GL3_1} в области $\Omega$ 

будем называть функцию

$u=u(x,y),$ удовлетворяющую уравнению \eqref{GL3_1} в области

$\Omega^{+}\cup\Omega^{-},$ и такую, что $y^{\,1-\alpha}u\in{C(\bar{\Omega})},$ 

$|x|^b u_{x}\in{C(\Omega)},$ $u_{xx},$ $D_{\,0y}^{\,\alpha} u\in{C(\Omega^{+}\cup\Omega^{-})}.$



\smallskip



{\sc\small Задача Коши.}

{\it Найти регулярное в области $\Omega$ решение

уравнения} \eqref{GL3_1}, {\it удовлетворяющее условию

\begin{equation}\label{UslSMCoshi}

\lim\limits_{y \rightarrow 0} D_{0y}^{\alpha-1}u(x,y)=\varphi(x), \quad  -\infty<x<\infty,

\end{equation}

где $\varphi(x)$ "--- заданная функция}.



\smallskip



В работе \cite{KhushtovaVestn2} доказана следующая теорема единственности.



\smallskip



\begin{newthm}{Теорема}  

Существует не более одного регулярного решения задачи Коши {\rm \eqref{GL3_1}, \eqref{UslSMCoshi} } в классе функций$,$ удовлетворяющих для некоторой положительной постоянной $k$ условию

\begin{equation} \label{AnalogTikhonovaInfty}

\lim\limits_{|x|\to \infty}y^{1-\alpha}\,u(x,y)\exp{\big(-k\,|x|^{\,\frac{2}{2-\alpha}} \big)}=0.

\end{equation}  

\end{newthm}



\smallskip





Отметим, что условие \eqref{AnalogTikhonovaInfty} имеет место и в случае задачи Коши для уравнения дробной диффузии.  

Теоремы единственности решения задачи Коши в~классе ограниченных функций  и~в~классе функций экспоненциального роста для многомерного уравнения дробной диффузии с~оператором Капуто были доказаны в работе А.~Н.~Кочубея \cite{Kochubei_2},

теорема единственности решения задачи Коши в классе функций экспоненциального роста для многомерного диффузионно"=волнового уравнения с~оператором Джрбашяна"--~Нерсесяна  "--- в~работе А.~В.~Псху~\cite{Pskhu2}. 

В этих же работах построены примеры, показывающие, что увеличение показателя $2/(2-\alpha)$ ведет к~утрате единственности решения соответствующих задач. 

При $\alpha=1$ условие \eqref{AnalogTikhonovaInfty} совпадает с~хорошо известным условием Тихонова 

$$

\lim\limits_{|x|\to \infty}u(x,y)\exp{\big(-k\,|x|^{2} \big)}=0.

$$

Как известно, пример неединственности решения задачи Коши для уравнения теплопроводности был впервые построен А.~Н.~Тихоновым в~основополагающей работе~\cite{tikhonovUsl-1} (см.~также \cite{tikhonovUsl-2}).



В данной работе эти идеи развиваются применительно к~задаче \eqref{GL3_1}, \eqref{UslSMCoshi}.



\smallskip



\Section{Неулучшаемость показателя степени в условии единственности решения задачи Коши}

Рассмотрим функцию

\begin{equation} \label{T(z,zeta)}  

T_{\mu}(z,\zeta)=\sum\limits_{n=0}^{\infty}\frac{z^{\,\beta+n}}{\,n!\,\Gamma(1+\beta+n)}\,\phi\left(-\rho,\mu-\alpha\beta-\alpha n; -\zeta\right),

\end{equation}    

где  $\phi\,(-\rho,\delta;z)$ "--- {\it функция Райта}, определяемая степенным рядом \cite{Pskhu, GorenfloLuchkoMainardi}

$$

\phi\,(-\rho,\delta;z)=\sum\limits_{n=0}^{\infty}\frac{z^{n}}{\,n!\,\Gamma(\delta-\rho\, n)}.

$$



Далее будем считать, что $\rho\in(0;1).$ 

Для любых $\nu$, $\delta\in \mathbb{R}$ справедливы формулы  \cite[с.~26, 44]{Pskhu}

\begin{gather}\label{diff_x}

\frac{d}{dz}\phi(-\rho,\delta;-z)=-\phi(-\rho,\delta-\rho;-z),

\\

\label{diff_Right}

D_{ay}^{\nu} \,y^{\delta-1}\phi(-\rho,\delta;-c\,y^{-\rho})=y^{\delta-\nu-1}\phi(-\rho,\delta-\nu;-c\,y^{-\rho}).

\end{gather}



\smallskip



\begin{lemma}[1]Если $\delta\geqslant 1,$ то для любого положительного $z$ справедливо нера\-вен\-ство~{\rm \cite[с.~47]{Pskhu} }

\begin{equation}\label{Right_delta>1}

\phi\,(-\rho,\delta;-z)\leqslant \frac{1}{\Gamma\left( \delta \right)}

\exp\left[-(1-\rho)\,\rho^{\frac{\rho}{1-\rho}}\, z^{\frac{1}{1-\rho}} \right].

\end{equation}

\end{lemma}



\smallskip



\begin{lemma}[2]Если $\delta<1,$ то для любых положительных $x$ и $y,$ 

$a\in(0;x),$ $\omega\in(1/2; \omega_0),$ $\omega_0=\min\{1,1/(2\rho)\},$

справедливы неравенства~{\rm \cite[с.~49]{Pskhu}}

\begin{gather}\label{Right-0}

\Bigl|\,y^{\,\delta-1}\phi\Bigl(-\rho,\delta;-\frac{x}{y^{\rho}}\Bigr)\Bigr|

\leqslant 

\frac{\Gamma\left( (1-\delta)/\rho\right)}{\,\rho\, \pi \left[aC_{\rho}(\rho,\omega) \right]^{(1-\delta)/\rho}}\,

\phi\Bigl(-\rho,1;-\frac{x-a}{y^{\rho}}\Bigr),

\\

\label{Right-2}

\Bigl|y^{\,\delta-1}\phi\Bigl(-\rho,\delta;-\frac{x}{y^{\rho}}\Bigr)\Bigr|

\leqslant \frac{\Gamma\left( 1-\delta\right)}{\pi \left[yC_{\rho}(1,\omega) \right]^{1-\delta}},

\end{gather}

где 

$C_{\rho}(\rho,\omega)=\frac1{\rho\pi}{\cos \rho\omega\pi},$ $C_{\rho}(1,\omega)=-\frac1{\pi}{\cos \omega\pi}.

$

\end{lemma}



\smallskip



Как следует из \eqref{Right-2}, ряд \eqref{T(z,zeta)} сходится абсолютно для любых  $z\in\mathbb{C},$ $\zeta>0,$ $\rho, \alpha, \beta\in(0;1),$ $\mu\in\mathbb{R}.$



\smallskip



Докажем следующую лемму.



\smallskip

\begin{lemma}[3]Если $0<\alpha<\rho<1,$ $\mu\geqslant 0,$ то при любых $x\in\mathbb{R},$ $y>0$ справедливо неравенство

\begin{equation} \label{T_muOtsenka}  

\Bigl| y^{\mu-1} T_{\mu}\Bigl(\frac{x^2}{4y^{\alpha}},\frac{1}{y^{\rho}}\Bigr)\Bigr|

\leqslant 

\frac{C

|x|^{\frac{2[1-\rho(1-\beta)]}{2\rho-\alpha}} y^{\mu}}{\Gamma(\mu+1)}\exp\left[ k_1  |x|^{\frac{2\rho}{2\rho-\alpha}}-k_2   y^{-\frac{\rho}{1-\rho}}\right],

\end{equation}    

где $C,$ $k_1,$ $k_2$ "--- положительные постоянные$,$ зависящие только от 

$\alpha,$ $\beta$ и $\rho.$

\end{lemma}



\smallskip



\begin{proof}

Из \eqref{Right-0} следует оценка

$$

\Bigl| y^{-\alpha\beta-\alpha n-1} \phi\Bigl(-\rho,-\alpha\beta-\alpha n; -\frac{1}{y^{\rho}}\Bigr)\Bigr|

\leqslant

 \frac{\Gamma\left( 1/\rho+\varepsilon\beta+\varepsilon n\right)}{ \rho \pi\left(aC_{\rho} \right)^{1/\rho+\varepsilon\beta+\varepsilon n}}

\phi\Bigl(-\rho,1; -\frac{1-a}{y^{\rho}}\Bigr),

$$

где $\varepsilon=\alpha/\rho,$ $a\in(0;1),$ $C_{\rho}=C_{\rho}(\rho,\omega).$



С помощью последней оценки получим

\begin{multline}

\Bigl| y^{\,\mu-1} T_{\mu}\Bigl(\frac{x^2}{4y^{\alpha}},\frac{1}{y^{\rho}}\Bigr)\Bigr|=

\\

=\left|

 D_{0y}^{-\mu}

  \sum\limits_{n=0}^{\infty}\frac{|x|^{2\beta+2n}}{ 2^{2\beta+2n}n! \Gamma(1+\beta+n)}

y^{-\alpha\beta-\alpha n-1}\phi \Bigl(-\rho,-\alpha\beta-\alpha n; -\frac{1}{y^{\rho}}\Bigr)

 \right|\leqslant

\\

 \label{WrightRyadZ}  

\leqslant

\frac{ |x|^{2\beta}y^{\mu}\phi\bigl(-\rho,\mu+1; -\frac{1-a}{y^{\rho}}\bigr)}{2^{2\beta}\rho \pi \left(aC_{\rho} \right)^{1/\rho+\varepsilon\beta}}

\sum\limits_{n=0}^{\infty}

\frac{\Gamma\left(1/\rho+\varepsilon\beta+\varepsilon n\right)}{n!

\Gamma(1+\beta+n)}

\Bigl( \frac{x^2}{4\left(aC_{\rho} \right)^{\varepsilon}}\Bigr)^{n}.

\end{multline}



Степенной ряд, стоящий справа от неравенства \eqref{WrightRyadZ}, представляет собой 

обобщенную функцию Райта

$$

{_1\Psi_{1}}\left[ \,z \,\bigg|

\begin{array}{ll}

 \big(1/\rho+\varepsilon\beta, \,\varepsilon\,\big)\\

\big(\,1+\beta, 1\,\big)\\

\end{array}

\right]=

\sum\limits_{n=0}^{\infty}

\frac{\Gamma\left(1/\rho+\varepsilon\beta+\varepsilon n\right)}{n!

\,\Gamma(1+\beta+n)} z^{n}, \; \; z=\frac{x^{\,2}}{4(aC_{\rho})^{\varepsilon}}.

$$

Из ее асимптотических свойств при $z\to \infty$ следует~\cite{Wright}

$$

{_1\Psi_{1}}\left[ \,z\,\bigg|

\begin{array}{ll}

 \big(1/\rho+\varepsilon\beta, \,\varepsilon\,\big)\\

\big(\,1+\beta, 1\,\big)\\

\end{array}

\right]=Z^{\,1/\rho-\beta(1-\varepsilon)-1}\,e^{Z}

\left[ \sum\limits_{m=0}^{M-1}A_m Z^{-m}+O\left(Z^{-M}\right) \right],

$$

где

$Z=(2-\varepsilon)\,\varepsilon^{\,\frac{\varepsilon}{2-\varepsilon}}z^{\,\frac{1}{2-\varepsilon}},$  $\varepsilon<2,$ а коэффициенты $A_m,$ $m=0, 1, 2, ...,$ зависят только от $\rho,$ $\varepsilon$ и $\beta.$



Учитывая, что $a$ "--- произвольное число из интервала $(0;1),$ из последней оценки и~соотношений \eqref{WrightRyadZ}, \eqref{Right_delta>1} получаем \eqref{T_muOtsenka}. \hfill \end{proof}



\smallskip



Покажем, что функция $y^{\mu-1} T_{\mu}\bigl(\frac{x^2}{4y^{\alpha}},\frac{1}{y^{\rho}}\bigr)$ является решением уравнения \eqref{GL3_1}. 

Из \eqref{T(z,zeta)}, \eqref{diff_x}  и \eqref{diff_Right} имеем

\begin{multline}

B_x T_{\mu}\Bigl(\frac{x^2}{4y^{\alpha}},\frac{1}{y^{\rho}}\Bigr)=

\\

=y^{-\alpha}\sum\limits_{n=1}^{\infty}\frac{1}{(n-1)! \Gamma(\beta+n)}

\Bigl(\frac{x^2}{4y^{\alpha}}\Bigr)^{\beta+n-1}

\phi\Bigl(-\rho,\mu-\alpha\beta-\alpha n; -\frac{1}{y^{\rho}}\Bigr)=

\\

=y^{-\alpha}\sum\limits_{k=0}^{\infty}\frac{1}{k! \Gamma(1+\beta+k)}

\Bigl(\frac{x^2}{4y^{\alpha}}\Bigr)^{\beta+k}

\phi\Bigl(-\rho,\mu-\alpha-\alpha\beta-\alpha k; -\frac{1}{y^{\rho}}\Bigr)=

\\ 

=y^{-\alpha} T_{\mu-\alpha}\Bigl(\frac{x^2}{\,4y^{\alpha}},\frac{1}{y^{\rho}}\Bigr),

\end{multline}   

\begin{equation}

\label{DT(z,zeta)}  

D_{0y}^{\nu} y^{\mu-1} T_{\mu}\Bigl(\frac{x^2}{4y^{\alpha}},\frac{1}{y^{\rho}}\Bigr)=

y^{\mu-\nu-1} T_{\mu-\nu}\Bigl(\frac{x^2}{4y^{\alpha}},\frac{1}{y^{\rho}}\Bigr).

\end{equation}



Из последних двух равенств следует

\begin{equation} \label{BDyT}  

\left(B_x-D_{0y}^{\alpha}\right)y^{\mu-1} T_{\mu}\Bigl(\frac{x^2}{4y^{\alpha}}, \frac{1}{y^{\rho}}\Bigr)=0.

\end{equation}   





Из равенства \eqref{DT(z,zeta)} и оценки \eqref{T_muOtsenka} также имеем

\begin{equation} \label{D(alpha-1)T}  

\lim\limits_{y \rightarrow 0} D_{0y}^{\alpha-1} y^{\mu-1} T_{\mu}\Bigl(\frac{x^2}{4y^{\alpha}}, \frac{1}{y^{\rho}}\Bigr)=0.

\end{equation}   



Таким образом, из соотношений  \eqref{T_muOtsenka}, \eqref{BDyT} и \eqref{D(alpha-1)T} следует, что нетривиальная функция 

$$

u(x,y)=y^{\mu-1} T_{\mu}\Bigl(\frac{x^2}{4y^{\alpha}},\frac{1}{y^{\rho}}\Bigr)

$$

является решением однородного уравнения \eqref{GL3_1}, удовлетворяет однородному условию \eqref{UslSMCoshi} $(\varphi(x)\equiv 0)$ и для любого положительного сколь угодно малого $\delta,$ 

некоторых положительных постоянных $k$ и $C,$ а также параметра 

$\rho,$  выбранного из условия 

$\rho=\frac{\alpha}{2}\frac{2+\delta(2-\alpha)}{\alpha+\delta(2-\alpha)},$

имеет место оценка

$$

\left| u(x,y) \right|\leqslant C \exp\left(k |x|^{\delta+\frac{2}{2-\alpha}} \right).

$$



Последняя оценка показывает, что увеличение показателя  $2/(2-\alpha)$ в~условии \eqref{AnalogTikhonovaInfty} ведет к~неединственности решения задачи  \eqref{GL3_1},  \eqref{UslSMCoshi} и~в~этом смысле условие  \eqref{AnalogTikhonovaInfty} является необходимым.





\smallskip









\Section[N]{Заключение}

В данной работе показывается, что показатель степени в~условии, обеспечивающем единственность решения задачи Коши для уравнения дробной диффузии с оператором Бесселя, является предельным, и его увеличение приводит к нарушению единственности. 



\smallskip





\AdditionalContent{Конкурирующие интересы}{Конкурирующих интересов не имею.} 



\AdditionalContent{Авторская ответственность}{Я несу полную ответственность за предоставление окончательной версии рукописи в печать. Окончательная версия рукописи мною одобрена.} 





\AdditionalContent{Финансирование}{Исследование выполнялось без финансирования.}



\smallskip





\begin{thebibliography}{99}



\RBibitem{Drob}

\by Нахушев~А.~М.

\book Дробное исчисление и его применение

\yr 2003

\publ Физматлит

\publaddr М.

\totalpages 271

%\sortdate 1/1/2003

\zmath{https://zbmath.org/?q=an:1066.26005}

%\misctransl

%\by Nakhushev~A.~M.

%\book Drobnoe ischislenie i ego primenenie \rm [Fractional calculus and its applications]

%\yr 2003

%\publ Fizmatlit

%\publaddr Moscow

%\lang In Russian

%\totalpages 271

%%\sortdate 1/1/2003



\RBibitem{Pskhu}

\by Псху~А.~В.

\book Уравнения в частных производных дробного порядка

\yr 2005

\publ Наука

\publaddr М.

\totalpages 199

%\sortdate 1/1/2005

\zmath{https://zbmath.org/?q=an:1193.35245}

%\misctransl

%\by Pskhu~A.~V.

%\book Uravneniia v chastnykh proizvodnykh drobnogo poriadka \rm [Partial differential equations of fractional order]

%\yr 2005

%\publ Nauka

%\publaddr Moscow

%\lang In Russian

%\totalpages 199

%%\sortdate 1/1/2005



\Bibitem{Alexiades}

\by Alexiades~V.

\paper Generalized axially symmetric heat potentials and singular parabolic initial boundary value problems

\jour Arch. Rational Mech. Anal.

\yr 1982

\vol 79

\issue 4

\pages 325--350

%\citnumscopus 3

%\sortdate 1/1/1982

\crossref{10.1007/BF00250797}

\mathscinet{http://www.ams.org/mathscinet-getitem?mr=656798}

\zmath{https://zbmath.org/?q=an:0511.35050}

\scopus{http://www.scopus.com/record/display.url?origin=inward&eid=2-s2.0-0020402046}



\Bibitem{Calton}

\by Calton~D.

\paper Cauchy's problem for a singular parabolic partial differential equation

\jour J.~Diff. Equations

\yr 1970

\vol 8

\issue 2

\pages 250--257

%\citnumscopus 8

%\sortdate 1/1/1970

\crossref{10.1016/0022-0396(70)90004-5}

\mathscinet{http://www.ams.org/mathscinet-getitem?mr=271547}

\scopus{http://www.scopus.com/record/display.url?origin=inward&eid=2-s2.0-49849107701}



\Bibitem{Kepinski_2}

\by K\c{e}pi\'nski~S.

\paper \"Uber die Differentialgleichung $\frac{\partial^2 z }{\partial x^2}+\frac{m+1}{x}\frac{\partial z }{\partial x} -\frac{n}{x}\frac{\partial z }{\partial t}=0$

\jour Math. Ann.

\yr 1905

\vol 61

\issue 3

\pages 397--405

\lang In German

%\sortdate 1/1/1905

\mathscinet{http://www.ams.org/mathscinet-getitem?mr=1511354}

\zmath{https://zbmath.org/?q=an:36.0416.01}



\RBibitem{Tersenov}

\by Терсенов~С.~А.

\book Параболические уравнения с~меняющимся направлением времени

\yr 1985

\publ Наука

\publaddr М.

\totalpages 105

%\sortdate 1/1/1985

%\misctransl

%\by Tersenov~S.~A.

%\book Parabolicheskie uravneniia s~meniaiushchimsia napravleniem vremeni \rm [Parabolic Equations with a Changing Time Direction]

%\yr 1985

%\publ Nauka

%\publaddr Moscow

%\totalpages 105

%%\sortdate 1/1/1985



\RBibitem{Matiychuk}

\by Матiйчук~M.~I.

\book Параболiчнi сингулярнi крайовi задачi

\yr 1999

\publ Iн-т математики НАН Укра\"\iни

\publaddr Ки\"\iв

\totalpages 176

%%\sortdate 1/1/1999

%\misctransl

%\by Matiichuk~M.~I.

%\book Parabolichni singuliarni kraiovi zadachi \rm [Parabolic Singular Boundary-Value Problems]

%\yr 1999

%\publ Institute of Mathematics of the Ukrainian National Academy of Sciences

%\publaddr Kiev

%\lang In Ukrainian

%\totalpages 176

%%\sortdate 1/1/1999



\RBibitem{Kipriyanov-Monographiya}

\by Киприянов~И.~А.

\book Сингулярные эллиптические краевые задачи

\yr 1997

\publ Наука

\publaddr М.

\totalpages 208

%\sortdate 1/1/1997

%\misctransl

%\by Kipriyanov~I.~A.

%\book Singuliarnye ellipticheskie kraevye zadachi \rm [Singular Elliptic Boundary-Value Problems]

%\yr 1997

%\publ Nauka

%\publaddr Moscow

%\lang In Russian

%\totalpages 208

%%\sortdate 1/1/1997



\RBibitem{Kipriyanov}

\by Киприянов~И.~А., Катрахов~В.~В., Ляпин~В.~М.

\paper О краевых задачах в областях общего вида для сингулярных параболических систем уравнений

\jour Докл. АН СССР

\yr 1976

\vol 230

\issue 6

\pages 1271--1274

%\sortdate 1/1/1976

\mathnet{http://mi.mathnet.ru/rus/dan40691}

\mathscinet{http://www.ams.org/mathscinet-getitem?mr=0425366}

\zmath{https://zbmath.org/?q=an:0354.35056}

%\misctransl

%\by Kipriyanov~I.~A., Katrakhov~V.~V., Lyapin~V.~M.

%\paper On boundary value problems in domains of general type for singular parabolic systems of equations

%\jour Sov. Math., Dokl.

%\yr 1976

%\vol 17

%\pages 1461--1464

%%\sortdate 1/1/1976

%\mathscinet{http://www.ams.org/mathscinet-getitem?mr=0425366}

%\zmath{https://zbmath.org/?q=an:0354.35056}



\RBibitem{Sitnik}

\by Ситник~С.~М.

\yr 2016

\publaddr Воронеж

\thesis Применение операторов преобразования Бушмана--Эрдейи и их обобщений в~теории дифференциальных уравнений с~особенностями в коэффициентах

\thesisinfo Дис.~\dots\ д-ра физ.-мат. наук: 01.01.02

\totalpages 307

%\sortdate 1/1/2016

%\misctransl

%\by Sitnik~S.~M.

%\yr 2016

%\publaddr Voronezh

%\thesis Application of the Bushman--Erd\'elyi transformation operators and their generalizations in the theory of differential equations with singularities in the coefficients

%\thesisinfo Dr. Phys. Math. Sci. Thesis

%\lang In Russian

%\totalpages 307

%%\sortdate 1/1/2016



\RBibitem{KatrakhovSitnik}

\by Катрахов~В.~В., Ситник~С.~М.

\paper Метод операторов преобразования и краевые задачи для сингулярных эллиптических уравнений

\inbook Сингулярные дифференциальные уравнения

\serial СМФН

\yr 2018

\vol 64

\issue 2

\pages 211--426

\publ Российский университет дружбы народов

\publaddr М.

%\sortdate 1/1/2018

\mathnet{http://mi.mathnet.ru/rus/cmfd355}

\crossref{10.22363/2413-3639-2018-64-2-211-426}

%\misctransl

%\by Katrakhov~V.~V., Sitnik~S.~M.

%\paper The transmutation method and boundary-value problems for singular elliptic equations

%\inbook Singular differential equations

%\serial CMFD

%\yr 2018

%\vol 64

%\issue 2

%\pages 211--426

%\publ Peoples' Friendship University of Russia

%\publaddr Moscow

%\lang In Russian

%%\sortdate 1/1/2018

%\crossref{https://doi.org/10.22363/2413-3639-2018-64-2-211-426}



\RBibitem{Kochubei_2}

\by Кочубей~А.~Н.

\paper Диффузия дробного порядка

\jour Дифференц. уравнения

\yr 1990

\vol 26

\issue 4

\pages 660--670

%\sortdate 1/1/1990

\mathnet{http://mi.mathnet.ru/rus/de7144}

\mathscinet{http://www.ams.org/mathscinet-getitem?mr=1061448}

\zmath{https://zbmath.org/?q=an:0712.35049}

%\transl

%\by Kochubei~A.~N.

%\paper Diffusion of fractional order

%\jour Differ. Equ.

%\yr 1990

%\vol 26

%\issue 4

%\pages 485--492

%%\sortdate 1/1/1990

%\mathscinet{http://www.ams.org/mathscinet-getitem?mr=1061448}

%\zmath{https://zbmath.org/?q=an:0729.35064}



\RBibitem{Kochubei_1}

\by Кочубей~А.~Н.

\paper Задача Коши для эволюционных уравнений дробного порядка

\jour Дифференц. уравнения

\yr 1989

\vol 25

\issue 8

\pages 1359--1368

%\sortdate 1/1/1989

\mathnet{http://mi.mathnet.ru/rus/de6938}

\mathscinet{http://www.ams.org/mathscinet-getitem?mr=1014153}

%\transl

%\by Kochubei~A.~N.

%\paper The Cauchy problem for evolution equations of fractional order

%\jour Differ. Equ.

%\yr 1989

%\vol 25

%\issue 8

%\pages 967--974

%%\sortdate 1/1/1989

%\mathscinet{http://www.ams.org/mathscinet-getitem?mr=1014153}



\RBibitem{Kochubei_3}

\by Кочубей~А.~Н., Эйдельман~С.~Д.

\paper Задача Коши для эволюционных уравнений дробного порядка

\jour Докл. РАН

\yr 2004

\vol 394

\issue 2

\pages 159--161

%\sortdate 1/1/2004

\mathnet{http://mi.mathnet.ru/rus/dan1624}

\mathscinet{http://www.ams.org/mathscinet-getitem?mr=2089744}

\zmath{https://zbmath.org/?q=an:1129.35427}

%\misctransl

%\by Kochubeii~A.~N., \'Eidel'man~S.~D.

%\paper The Cauchy problem for evolution equations of fractional order

%\jour Dokl. Akad. Nauk

%\yr 2004

%\vol 394

%\issue 2

%\pages 159--161

%\lang In Russian

%%\sortdate 1/1/2004



\Bibitem{Kochubei_4}

\by Kochubei~A.~N.

\paper Asymptotic Properties of Solutions of the Fractional Diffusion-Wave Equation

\jour Fract. Calc. Appl. Anal.

\yr 2014

\vol 17

\issue 3

\pages 881--896

%\citnumscopus 9

%\sortdate 1/1/2014

\crossref{10.2478/s13540-014-0203-3}

\mathscinet{http://www.ams.org/mathscinet-getitem?mr=3260311}

\zmath{https://zbmath.org/?q=an:1323.35197}

\elib{http://elibrary.ru/item.asp?id=24541812}

\scopus{http://www.scopus.com/record/display.url?origin=inward&eid=2-s2.0-84904312827}



\Bibitem{Kochubei_5}

\by Kochubei~A.~N.

\paper Cauchy problem for fractional diffusion-wave equations with variable coefficients

\jour Applicable Analysis

\yr 2014

\vol 93

\issue 10

\pages 2211--2242

%\citnumscopus 7

%\sortdate 1/1/2014

\crossref{10.1080/00036811.2013.875162}

\mathscinet{http://www.ams.org/mathscinet-getitem?mr=3240385}

\zmath{https://zbmath.org/?q=an:1297.35274}

\elib{http://elibrary.ru/item.asp?id=24642244}

\scopus{http://www.scopus.com/record/display.url?origin=inward&eid=2-s2.0-84905387477}



\RBibitem{Pskhu2}

\by Псху~А.~В.

\paper Фундаментальное решение диффузионно-волнового уравнения дробного порядка

\jour Изв. РАН. Сер. матем.

\yr 2009

\vol 73

\issue 2

\pages 141--182

%\sortdate 1/1/2009

\mathnet{http://mi.mathnet.ru/rus/im2429}

\crossref{10.4213/im2429}

\mathscinet{http://www.ams.org/mathscinet-getitem?mr=2532450}

\zmath{https://zbmath.org/?q=an:1172.26001}

\adsnasa{http://adsabs.harvard.edu/cgi-bin/bib_query?2009IzMat..73..351P}

\elib{http://elibrary.ru/item.asp?id=20425204}

%\transl

%\by Pskhu~A.~V.

%\paper The fundamental solution of a diffusion-wave equation of fractional order

%\jour Izv. Math.

%\yr 2009

%\vol 73

%\issue 2

%\pages 351--392

%%\sortdate 1/1/2009

%\crossref{https://doi.org/10.1070/IM2009v073n02ABEH002450}

%\mathscinet{http://www.ams.org/mathscinet-getitem?mr=2532450}

%\zmath{https://zbmath.org/?q=an:1172.26001}

%\isi{http://gateway.isiknowledge.com/gateway/Gateway.cgi?GWVersion=2&SrcApp=PARTNER_APP&SrcAuth=LinksAMR&DestLinkType=FullRecord&DestApp=ALL_WOS&KeyUT=000266177900006}

%\elib{http://elibrary.ru/item.asp?id=13597908}

%\scopus{http://www.scopus.com/record/display.url?origin=inward&eid=2-s2.0-65349154473}



\RBibitem{Pskhu-SibirJournal}

\by Псху~А.~В.

\paper Уравнение дробной диффузии с оператором дискретно распределенного дифференцирования

\jour Сиб. электрон. матем. изв.

\yr 2016

\vol 13

\pages 1078--1098

%\sortdate 1/1/2016

\mathnet{http://mi.mathnet.ru/rus/semr736}

\crossref{10.17377/semi.2016.13.086}

\mathscinet{http://www.ams.org/mathscinet-getitem?mr=3580054}

\zmath{https://zbmath.org/?q=an:1370.35268}

\elib{http://elibrary.ru/item.asp?id=28127183}

%\misctransl

%\by Pskhu~A.~V.

%\paper Fractional diffusion equation with discretely distributed differentiation operator

%\jour Sib. \`Elektron. Mat. Izv.

%\yr 2016

%\vol 13

%\pages 1078--1098

%\lang In Russian

%%\sortdate 1/1/2016

%\crossref{https://doi.org/10.17377/semi.2016.13.086}



\Bibitem{Pskhu3}

\by Pskhu~A.~V.

\paper Multi-time fractional diffusion equation

\jour Eur. Phys. J.~Spec. Top.

\yr 2013

\vol 222

\issue 8

\pages 1939--1950

%\citnumscopus 2

%\sortdate 1/1/2013

\crossref{10.1140/epjst/e2013-01975-y}

\scopus{http://www.scopus.com/record/display.url?origin=inward&eid=2-s2.0-84885034281}



\RBibitem{Pskhu4}

\by Псху~А.~В.

\paper Первая краевая задача для дробного диффузионно-волнового уравнения в~нецилиндрической области

\jour Изв. РАН. Сер. матем.

\yr 2017

\vol 81

\issue 6

\pages 158--179

%\sortdate 1/1/2017

\mathnet{http://mi.mathnet.ru/rus/im8520}

\crossref{10.4213/im8520}

\mathscinet{http://www.ams.org/mathscinet-getitem?mr=3733319}

\zmath{https://zbmath.org/?q=an:06844352}

\adsnasa{http://adsabs.harvard.edu/cgi-bin/bib_query?2017IzMat..81.1212P}

\elib{http://elibrary.ru/item.asp?id=30737847}

%\transl

%\by Pskhu~A.~V.

%\paper The first boundary-value problem for a fractional diffusion-wave equation in a non-cylindrical domain

%\jour Izv. Math.

%\yr 2017

%\vol 81

%\issue 6

%\pages 1212--1233

%%\sortdate 1/1/2017

%\crossref{https://doi.org/10.1070/IM8520}

%\mathscinet{http://www.ams.org/mathscinet-getitem?mr=3733319}

%\zmath{https://zbmath.org/?q=an:06844352}

%\isi{http://gateway.isiknowledge.com/gateway/Gateway.cgi?GWVersion=2&SrcApp=PARTNER_APP&SrcAuth=LinksAMR&DestLinkType=FullRecord&DestApp=ALL_WOS&KeyUT=000418891300007}

%\scopus{http://www.scopus.com/record/display.url?origin=inward&eid=2-s2.0-85040983678}



\RBibitem{VorKilbas}

\by Ворошилов~А.~А., Килбас~А.~А.

\paper Задача типа Коши для диффузионно-волнового уравнения с частной производной Римана--Лиувилля

\jour Докл. РАН

\yr 2006

\vol 406

\issue 1

\pages 12--16

%\sortdate 1/1/2006

\mathnet{http://mi.mathnet.ru/rus/dan1045}

\mathscinet{http://www.ams.org/mathscinet-getitem?mr=2256848}

\zmath{https://zbmath.org/?q=an:1155.35396}

\elib{http://elibrary.ru/item.asp?id=9176502}

%\misctransl

%\by Voroshilov~A.~A., Kilbas~A.~A.

%\paper A Cauchy-type problem for a diffusion-wave equation with a Riemann-Liouville partial derivative

%\jour Dokl. Akad. Nauk

%\yr 2006

%\vol 406

%\issue 1

%\pages 12--16

%\lang In Russian

%%\sortdate 1/1/2006



\RBibitem{VorKilbas2}

\by Ворошилов~А.~А., Килбас~А.~А.

\paper Задача Коши для диффузионно-волнового уравнения с~частной производной Капуто

\jour Дифференц. уравнения

\yr 2006

\vol 42

\issue 5

\pages 599--609

%\sortdate 1/1/2006

\mathnet{http://mi.mathnet.ru/rus/de11489}

\mathscinet{http://www.ams.org/mathscinet-getitem?mr=2292158}

\zmath{https://zbmath.org/?q=an:1123.35302}

\elib{http://elibrary.ru/item.asp?id=9245260}

%\transl

%\by Voroshilov~A.~A., Kilbas~A.~A.

%\paper The Cauchy problem for the diffusion-wave equation with the Caputo partial derivative

%\jour Differ. Equ.

%\yr 2006

%\vol 42

%\issue 5

%\pages 638--649

%%\citnumscopus 11

%%\sortdate 1/1/2006

%\crossref{https://doi.org/10.1134/S0012266106050041}

%\mathscinet{http://www.ams.org/mathscinet-getitem?mr=2292158}

%\zmath{https://zbmath.org/?q=an:1123.35302}

%\elib{http://elibrary.ru/item.asp?id=27786841}

%\scopus{http://www.scopus.com/record/display.url?origin=inward&eid=2-s2.0-33745304762}



\RBibitem{Gekk}

\by Геккиева~С.~Х.

\paper Задача Коши для обобщенного уравнения переноса с дробной по времени производной

\jour Доклады Адыгской (Черкесской) Международной академии наук

\yr 2000

\vol 5

\issue 1

\pages 16--19

%\sortdate 1/1/2000

%\misctransl

%\by Gekkieva~S.~Kh.

%\paper The Cauchy problem for the generalized transport equation with time-fractional derivative

%\jour Dokl. Adyg. (Cherkess.) Mezdunar. Akad. Nauk

%\yr 2000

%\vol 5

%\issue 1

%\pages 16--19

%\lang In Russian

%%\sortdate 1/1/2000



\RBibitem{Gekkieva1}

\by Геккиева~С.~Х.

\paper Краевая задача для обобщенного уравнения переноса с дробной производной в полубесконечной области

\jour Изв. КБНЦ РАН

\yr 2002

\issue 1(8)

\pages 6--8

%\sortdate 1/1/2002

%\misctransl

%\by Gekkieva~S.~Kh.

%\paper A boundary value problem for the generalized transfer equation with a fractional derivative in a semi-infinite domain

%\jour Izv. KBNTs RAN

%\yr 2002

%\issue 1(8)

%\pages 6--8

%\lang In Russian

%%\sortdate 1/1/2002



\Bibitem{Mainardi-1}

\by Mainardi~F.

\paper The fundamental solutions for the fractional diffusion-wave equation

\jour Appl. Math. Letters

\yr 1996

\vol 6

\issue 1

\pages 23--28

%\citnumscopus 343

%\sortdate 1/1/1996

\crossref{10.1016/0893-9659(96)00089-4}

\mathscinet{http://www.ams.org/mathscinet-getitem?mr=1419811}

\zmath{https://zbmath.org/?q=an:0879.35036}

\scopus{http://www.scopus.com/record/display.url?origin=inward&eid=2-s2.0-30244460855}



\Bibitem{Mainardi-3}

\by Mainardi~F.

\paper The time fractional diffusion-wave equation

\jour Radiophys. Quantum Electron.

\yr 1995

\vol 38

\issue 1-2

\pages 13--24

%\citnumscopus 39

%\sortdate 1/1/1995

\crossref{10.1007/BF01051854}

\mathscinet{http://www.ams.org/mathscinet-getitem?mr=1427165}

\scopus{http://www.scopus.com/record/display.url?origin=inward&eid=2-s2.0-0029427754}



\Bibitem{Mainardi-2}

\by Mainardi~F., Luchko~Yu., Pagnini~G.

\paper The fundamental solution of the space-time fractional diffusion equation

\jour Fract. Calc. Appl. Anal.

\yr 2001

\vol 4

\issue 2

\pages 153--192

\arxiv \href{https://arxiv.org/abs/cond-mat/0702419}{cond-mat/0702419} [cond-mat.stat-mech]

%\sortdate 1/1/2001

\mathscinet{http://www.ams.org/mathscinet-getitem?mr=1829592}

\zmath{https://zbmath.org/?q=an:1054.35156}



\Bibitem{Pagnini_1}

\by Pagnini~G.

\paper The M-Wright function as a generalization of the Gaussian density for fractional diffusion processes

\jour Fract. Calc. Appl. Anal.

\yr 2013

\vol 16

\issue 2

\pages 436--453

%\citnumscopus 22

%\sortdate 1/1/2013

\crossref{10.2478/s13540-013-0027-6}

\mathscinet{http://www.ams.org/mathscinet-getitem?mr=3033698}

\zmath{https://zbmath.org/?q=an:1312.33061}

\elib{http://elibrary.ru/item.asp?id=28682912}

\scopus{http://www.scopus.com/record/display.url?origin=inward&eid=2-s2.0-84879620364}



\Bibitem{Pagnini_Paradisi}

\by Pagnini~G., Paradisi~P.

\paper A stochastic solution with Gaussian stationary increments of the symmetric space-time fractional diffusion equation

\jour Fract. Calc. Appl. Anal.

\yr 2016

\vol 19

\issue 2

\pages 408--440

%\citnumscopus 9

%\sortdate 1/1/2016

\crossref{10.1515/fca-2016-0022}

\mathscinet{http://www.ams.org/mathscinet-getitem?mr=3513003}

\zmath{https://zbmath.org/?q=an:1341.60073}

\elib{http://elibrary.ru/item.asp?id=27104131}

\scopus{http://www.scopus.com/record/display.url?origin=inward&eid=2-s2.0-84969800118}



\Bibitem{Podlubny}

\by Podlubny~I.

\book Fractional differential equations. An introduction to fractional derivatives, fractional differential equations, to methods of their solution and some of their applications

\serial Mathematics in Science and Engineering

\yr 1999

\vol 198

\publ Academic Press

\publaddr San Diego, CA

\totalpages xxiv+340

%\sortdate 1/1/1999

\mathscinet{http://www.ams.org/mathscinet-getitem?mr=1658022}

\zmath{https://zbmath.org/?q=an:0924.34008}



\Bibitem{MathaiSaxenaHaubold}

\by Mathai~A.~M., Saxena~R.~K., Haubold~H.~J.

\book The $H$-function. Theory and Applications

\yr 2010

\publ Springer

\publaddr Dordrecht

\totalpages xiv+268~pp \nofrills

%\citnumscopus 411

%\sortdate 1/1/2010

\crossref{10.1007/978-1-4419-0916-9}

\mathscinet{http://www.ams.org/mathscinet-getitem?mr=2562766}

\zmath{https://zbmath.org/?q=an:1181.33001}

\scopus{http://www.scopus.com/record/display.url?origin=inward&eid=2-s2.0-84891408876}



\RBibitem{KhushtovaVestn2}

\by Хуштова~Ф.~Г.

\paper Задача Коши для уравнения параболического типа~с~оператором Бесселя и~частной производной~Римана--Лиувилля

\jour Вестн. Сам. гос. техн. ун-та. Сер. Физ.-мат. науки

\yr 2016

\vol 20

\issue 1

\pages 74--84

%\sortdate 1/1/2016

\mathnet{http://mi.mathnet.ru/rus/vsgtu1455}

\crossref{10.14498/vsgtu1455}

\zmath{https://zbmath.org/?q=an:06964474}

\elib{http://elibrary.ru/item.asp?id=26898117}

%\misctransl

%\by Khushtova~F.~G.

%\paper Cauchy problem for a parabolic equation with Bessel operator and Riemann--Liouville partial derivative

%\jour Vestn. Samar. Gos. Tekhn. Univ., Ser. Fiz.-Mat. Nauki \rm [J. Samara State Tech. Univ., Ser. Phys. Math. Sci.]

%\yr 2016

%\vol 20

%\issue 1

%\pages 74--84

%%\sortdate 1/1/2016



\Bibitem{GorenfloLuchkoMainardi}

\by Gorenflo~R., Luchko~Y., Mainardi~F.

\paper Analytical properties and applications of the Wright function

\jour Fract. Calc. Appl. Anal.

\yr 1999

\vol 2

\issue 4

\pages 383--414

\arxiv \href{https://arxiv.org/abs/math-ph/0701069}{math-ph/0701069}

%\sortdate 1/1/1999

\mathscinet{http://www.ams.org/mathscinet-getitem?mr=1752379}

\zmath{https://zbmath.org/?q=an:1027.33006}



\Bibitem{tikhonovUsl-1}

\by Tychonoff~A.

\paper Th\'eor\`emes d'unicit\'e pour l'\'equation de la chaleur

\jour Mat. Sb.

\yr 1935

\vol 42

\issue 2

\pages 199--216

\lang In French

%\sortdate 1/1/1935

\mathnet{http://mi.mathnet.ru/rus/sm6410}

\zmath{https://zbmath.org/?q=an:0012.35501|61.1203.05}



\RBibitem{tikhonovUsl-2}

\by Тихонов~А.~Н.

\paper Теоремы единственности для уравнения теплопроводности

\inbook Собрание научных трудов: в десяти томах

\yr 2009

\pages 371--382

\publ Наука

\publaddr М.

\volinfo Т.~II. Математика. Ч.~2. Вычислительная математика 1956--1979. Математическая физика 1933--1948. Ред.-сост. Т.~А.~Сушкевич, А.~В.~Гулин

%\sortdate 1/1/2009

\zmath{https://zbmath.org/?q=an:1200.01055}

%\misctransl

%\by Tikhonov~A.~N.

%\paper Uniqueness theorems for the heat equation

%\inbook Collection of scientific works. In ten volumes

%\yr 2009

%\pages 371--382

%\publ Nauka

%\publaddr Moscow

%\lang In Russian and French

%\volinfo Vol. 2: Mathematics. Part 2: Numerical mathematics 1956–1979. Mathematical physics 1933–1948. Edited by T.~A.~Sushkevich and A.~V.~Gulin

%%\sortdate 1/1/2009

%\zmath{https://zbmath.org/?q=an:1200.01055}



\Bibitem{Wright}

\by Wright~E.~M.

\paper The asymptotic expansion of the generalized hypergeometric function

\jour Proc. Lond. Math. Soc., II. Ser.

\yr 1940

\vol 46

\issue 1

\pages 389--408

%\citnumscopus 136

%\sortdate 1/1/1940

\crossref{10.1112/plms/s2-46.1.389}

\mathscinet{http://www.ams.org/mathscinet-getitem?mr=3876}

\zmath{https://zbmath.org/?q=an:0025.40402|66.1243.01}

\scopus{http://www.scopus.com/record/display.url?origin=inward&eid=2-s2.0-84963078619}



\end{thebibliography}



	





\newpage



\makeentitle



\vspace{-6mm}



%\AdditionalContent{Competing interests}{Specify what conflicts of interest are available.} 

\AdditionalContent{Competing interests}{I have no competing interests.} 

%\AdditionalContent{Competing interests}{We have ... (or) no competing interests.} 

%\AdditionalContent{Authors' contributions and responsibilities}{What is the author's responsibility for each author?}

\AdditionalContent{Author's Responsibilities}{I take full responsibility for submitting the final ma\-nuscript in print. I approved the final version of the manuscript.} 

%\AdditionalContent{Authors' contributions and responsibilities}{Each author has participated in the article concept development and in the manuscript writing. The authors are absolutely responsible for submitting the final manuscript in print. Each author has approved the final version of manuscript.} 



%\AdditionalContent{Funding}{Indicate the source(s) of your funding.}



\AdditionalContent{Funding}{This research received no specific grant from any funding agency in the public,

commercial, or not-for-profit sectors.}



\newpage



\begin{thebibliography}{10}



\Bibitem{Drob:2}

\lang In Russian

\by Nakhushev~A.~M.

\book Drobnoe ischislenie i ego primenenie \rm [Fractional calculus and its applications]

\publaddr Moscow

\publ Fizmatlit 

\totalpages 271

\yr 2003





\Bibitem{Pskhu:w}

\lang In Russian

\by Pskhu~A.~V.

\book Uravneniia v chastnykh proizvodnykh drobnogo poriadka \rm [Partial differential equations of fractional order]

\publaddr Moscow

\publ Nauka 

\totalpages 199 

\yr 2005







\Bibitem{Alexiades:w}

\by Alexiades~V.

\paper Generalized axially symmetric heat potentials and singular parabolic initial boundary value problems

\jour  Arch. Rational Mech. Anal. 

\yr 1982

\vol 79

\issue 4

\pages 325--350

\crossref{10.1007/BF00250797}







\Bibitem{Calton:w}

\by Calton~D.

\paper Cauchy's problem for a singular parabolic partial differential equation

\jour J.~Diff. Equations

\yr 1970

\vol 8

\issue 2

\pages 250--257

\crossref{10.1016/0022-0396(70)90004-5}







\Bibitem{Kepinski_2:w}

\by K\c{e}pi\'nski~S.

\paper \"Uber die Differentialgleichung $\frac{\partial^2 z }{\partial x^2}+\frac{m+1}{x}\frac{\partial z }{\partial x} -\frac{n}{x}\frac{\partial z }{\partial t}=0$

\jour   Math. Ann.

\yr 1905

\vol 61

\issue 3

\pages 397--405

\zmath{36.0416.01}

\lang In German



\Bibitem{Tersenov:en}

\lang In Russian

\by Tersenov~S.~A.

\book Parabolicheskie uravneniia s~meniaiushchimsia napravleniem vremeni \rm [Parabolic Equations with a Changing Time Direction]

\yr 1985

\publ Nauka

\publaddr Moscow

\totalpages 105



\Bibitem{Matiychuk:en}

\by Matiichuk~M.~I.

\book  Parabolichni singuliarni kraiovi zadachi \rm [Parabolic Singular Boundary-Value Problems]

\yr 1999

\publ Institute of Mathematics of the Ukrainian National  Academy of Sciences

\publaddr Kiev

\totalpages 176

\lang In Ukrainian





\Bibitem{Kipriyanov-Monographiya:en}

\lang In Russian

\by Kipriyanov~I.~A.

\book  Singuliarnye ellipticheskie kraevye zadachi \rm [Singular Elliptic Boundary-Value Problems]

\yr 1997

\publ Nauka

\publaddr Moscow

\totalpages 208



\Bibitem{Kipriyanov:en}

\by Kipriyanov~I.~A., Katrakhov~V.~V., Lyapin~V.~M.

\paper On boundary value problems in domains of general type for singular parabolic systems of equations

\jour Sov. Math., Dokl. 

\vol 17

\yr 1976

\pages 1461--1464 

\mathscinet{http://www.ams.org/mathscinet-getitem?mr=0425366}

\zmath{https://zbmath.org/?q=an:0354.35056}





\Bibitem{Sitnik:en}

\lang In Russian

\by Sitnik~S.~M. 

\thesis Application of the Bushman--Erd\'elyi transformation operators and their generalizations in the theory of differential equations with singularities in the coefficients

\thesisinfo Dr. Phys. Math. Sci. Thesis

\yr 2016

\publaddr Voronezh

\totalpages 307







\Bibitem{KatrakhovSitnik:en}

\lang In Russian

\by Katrakhov~V.~V., Sitnik~S.~M.

\paper The transmutation method and boundary-value problems for singular elliptic equations

\inbook Singular differential equations

\serial CMFD

\yr 2018

\vol 64

\issue 2

\pages 211--426

\publ Peoples' Friendship University of Russia

\publaddr Moscow

\crossref{10.22363/2413-3639-2018-64-2-211-426}







\Bibitem{Kochubei_2:en}

\by Kochubei~A.~N.

\paper Diffusion of fractional order

\jour Differ. Equ.

\yr 1990

\vol 26

\issue 4

\pages 485--492







\Bibitem{Kochubei_1:w}

\by Kochubei~A.~N.

\paper The Cauchy problem for evolution equations of fractional order

\jour Differ. Equ.

\yr 1989

\vol 25

\issue 8

\pages 967--974



\Bibitem{Kochubei_3:e}

\lang In Russian

\by Kochubeii~A.~N., \'Eidel'man~S.~D.

\paper The Cauchy problem for evolution equations of fractional order

\jour Dokl. Akad. Nauk

\yr 2004

\vol 394

\issue 2

\pages 159--161







\Bibitem{Kochubei_4:e}

\by Kochubei~A.~N.

\paper Asymptotic Properties of Solutions of the Fractional Diffusion-Wave Equation

\jour Fract. Calc. Appl. Anal.

\yr 2014

\vol 17

\issue 3

\pages 881--896

\crossref{10.2478/s13540-014-0203-3}





\Bibitem{Kochubei_5:e}

\by Kochubei~A.~N.

\paper Cauchy problem for fractional diffusion-wave equations with variable coefficients

\jour  Applicable Analysis

\yr 2014

\vol 93

\issue 10

\pages 2211--2242

\crossref{10.1080/00036811.2013.875162}







\Bibitem{Pskhu2:en}

\paper The fundamental solution of a diffusion-wave equation of fractional order

\by Pskhu~A.~V.

\jour Izv. Math.

\yr 2009

\vol 73

\issue 2

\pages 351--392

\crossref{10.1070/IM2009v073n02ABEH002450}

\isi{http://gateway.isiknowledge.com/gateway/Gateway.cgi?GWVersion=2&SrcApp=PARTNER_APP&SrcAuth=LinksAMR&DestLinkType=FullRecord&DestApp=ALL_WOS&KeyUT=000266177900006}

\elib{http://elibrary.ru/item.asp?id=13597908}

\scopus{http://www.scopus.com/record/display.url?origin=inward&eid=2-s2.0-65349154473}





\Bibitem{Pskhu-SibirJournal:en}

\lang In Russian

\by Pskhu~A.~V.

\paper Fractional diffusion equation with discretely distributed differentiation operator

\jour Sib. \`Elektron. Mat. Izv.

\yr 2016

\vol 13

\pages 1078--1098

\crossref{10.17377/semi.2016.13.086}







\Bibitem{Pskhu3:en}

\by Pskhu~A.~V.

\paper Multi-time fractional diffusion equation

\jour Eur. Phys. J.~Spec. Top.

\yr 2013

\vol 222

\issue 8

\pages 1939--1950

\crossref{10.1140/epjst/e2013-01975-y}





\Bibitem{Pskhu4:en}

\paper The first boundary-value problem for a fractional diffusion-wave equation in a non-cylindrical domain

\by Pskhu~A.~V.

\jour Izv. Math.

\yr 2017

\vol 81

\issue 6

\pages 1212--1233

\crossref{10.1070/IM8520}

\isi{http://gateway.isiknowledge.com/gateway/Gateway.cgi?GWVersion=2&SrcApp=PARTNER_APP&SrcAuth=LinksAMR&DestLinkType=FullRecord&DestApp=ALL_WOS&KeyUT=000418891300007}

\scopus{http://www.scopus.com/record/display.url?origin=inward&eid=2-s2.0-85040983678}







\Bibitem{VorKilbas:a} 

\lang In Russian

\by Voroshilov~A.~A., Kilbas~A.~A.

\paper A Cauchy-type problem for a diffusion-wave equation with a

Riemann-Liouville partial derivative

\jour Dokl. Akad. Nauk

\yr 2006

\vol 406

\issue 1

\pages 12--16



\Bibitem{VorKilbas2:en} 

\by Voroshilov~A.~A., Kilbas~A.~A.

\paper The Cauchy problem for the diffusion-wave equation with the Caputo partial derivative

\mathscinet{http://www.ams.org/mathscinet-getitem?mr=2292158}

\jour Differ. Equ.

\yr 2006

\vol 42

\issue 5

\pages 638--649

\crossref{10.1134/S0012266106050041}



\Bibitem{Gekk:en} 

\lang In Russian

\by Gekkieva~S.~Kh.

\paper The Cauchy problem for the generalized transport equation with time-fractional derivative

\jour Dokl. Adyg. (Cherkess.) Mezdunar. Akad. Nauk

\yr 2000

\vol 5

\issue 1

\pages 16--19







\Bibitem{Gekkieva1:en}

\lang In Russian

\by Gekkieva~S.~Kh.

\paper A boundary value problem for the generalized transfer equation with a fractional derivative in a semi-infinite domain

\jour Izv. KBNTs RAN

\yr 2002

\issue 1(8)

\pages 6--8











  





\Bibitem{Mainardi-1:en}

\by Mainardi~F.

\paper The fundamental solutions for the fractional diffusion-wave equation

\jour Appl.  Math. Letters

\yr 1996

\vol 6

\issue 1

\pages 23--28

\crossref{10.1016/0893-9659(96)00089-4}



\Bibitem{Mainardi-3}

\by Mainardi~F.

\paper  The time fractional diffusion-wave equation

\jour Radiophys. Quantum Electron.

\yr 1995

\vol 38

\issue 1-2

\pages 13--24

\crossref{10.1007/BF01051854}



\Bibitem{Mainardi-2:en}

\by Mainardi~F., Luchko~Yu., Pagnini~G.

\paper The fundamental solution of the space-time fractional diffusion equation

\jour Fract. Calc. Appl. Anal.

\yr 2001

\vol 4

\issue 2

\pages 153--192

\arxiv  	\href{https://arxiv.org/abs/cond-mat/0702419}{cond-mat/0702419} [cond-mat.stat-mech]















\Bibitem{Pagnini_1:en}

\by Pagnini~G.

\paper The M-Wright function as a generalization of the Gaussian density for fractional diffusion processes

\jour Fract. Calc. Appl. Anal.

\yr 2013

\vol 16

\issue 2

\pages 436--453

\crossref{10.2478/s13540-013-0027-6}





\Bibitem{Pagnini_Paradisi:en}

\by Pagnini~G., Paradisi~P.

\paper A stochastic solution with Gaussian stationary increments of the symmetric space-time fractional diffusion equation

\jour Fract. Calc. Appl. Anal.

\yr 2016

\vol 19

\issue 2

\pages 408--440

\crossref{10.1515/fca-2016-0022}





\Bibitem{Podlubny:en}

\by Podlubny~I.

\book Fractional differential equations. An introduction to fractional derivatives, fractional differential equations, to methods of their solution and some of their applications

\yr 1999

\publ Academic Press

\publaddr  San Diego, CA

\zmath{0924.34008}

\serial Mathematics in Science and Engineering

\vol 198

\totalpages xxiv+340





\Bibitem{MathaiSaxenaHaubold:en}

\by Mathai~A.~M., Saxena~R.~K., Haubold~H.~J.

\book The $H$-function. Theory and Applications

\yr 2010

\publ Springer

\publaddr  Dordrecht

\totalpages xiv+268~pp \nofrills

\crossref{10.1007/978-1-4419-0916-9}

\zmath{1181.33001}









\Bibitem{KhushtovaVestn2}

\by Khushtova~F.~G.

\paper Cauchy problem for a parabolic equation with Bessel operator and Riemann--Liouville partial derivative

\jour Vestn. Samar. Gos. Tekhn. Univ., Ser. Fiz.-Mat. Nauki \rm [J. Samara State Tech. Univ., Ser. Phys. Math. Sci.]

\yr 2016

\vol 20

\issue 1

\pages 74--84

\crossref{10.14498/vsgtu1455}





\Bibitem{GorenfloLuchkoMainardi:en}

\by Gorenflo~R., Luchko~Y., Mainardi~F.

\paper Analytical properties and applications of the Wright  function

\jour Fract. Calc. Appl. Anal.

\yr 1999

\vol 2

\issue 4

\pages 383--414

\arxiv \href{https://arxiv.org/abs/math-ph/0701069}{math-ph/0701069}







\Bibitem{tikhonovUsl-1:en}

\by Tychonoff~A.

\paper Th\'eor\`emes d'unicit\'e pour l'\'equation de la chaleur

\jour Mat. Sb.

\yr 1935

\vol 42

\issue 2

\pages 199--216

\lang In French

\mathnet{http://mi.mathnet.ru/msb6410}

\zmath{https://zbmath.org/?q=an:0012.35501|61.1203.05}









\Bibitem{tikhonovUsl-2:en}

\by Tikhonov~A.~N.

\paper Uniqueness theorems for the heat equation

\inbook Collection of scientific works. In ten volumes

\volinfo Vol. 2: Mathematics. Part 2: Numerical mathematics 1956–1979. Mathematical physics 1933–1948. Edited by T.~A.~Sushkevich and  A.~V.~Gulin

\yr 2009

\publ Nauka

\publaddr Moscow

\pages 371--382

\zmath{1200.01055}

\lang In Russian and French





\Bibitem{Wright:en}

\by Wright~E.~M.

\paper The asymptotic expansion of the generalized hypergeometric function

\jour Proc. Lond. Math. Soc., II. Ser. 

\yr 1940

\vol 46

\issue 1

\pages 389--408

\crossref{10.1112/plms/s2-46.1.389}

\zmath{0025.40402|66.1243.01}











\end{thebibliography}



%\end{document}

